\documentclass{article}
\usepackage[a4paper,top=3cm,bottom=3cm,left=2.5cm,right=2.5cm,marginparwidth=1.75cm]{geometry}
\title{Piano di Progetto}
\usepackage[T1]{fontenc}
%\usepackage{lmodern}
\usepackage[utf8]{inputenc}
\usepackage[italian]{babel}
\usepackage{amsmath}
\usepackage{booktabs}
\usepackage{longtable}
\usepackage{wrapfig}
\usepackage{mdframed}
\usepackage{hyperref}
\usepackage{adjustbox}
\usepackage{lastpage}

%% HEADER
\usepackage[sfdefault,lf]{carlito}
%% The 'lf' option for lining figures
%% The 'sfdefault' option to make the base font sans serif
\usepackage[parfill]{parskip}
\renewcommand*\oldstylenums[1]{\carlitoOsF #1}
\usepackage{fancyhdr}
%\usepackage{natbib}
%\usepackage{authblk}
\setlength{\headheight}{41pt}

\fancyhead[L]{WatchDoge: Piano di Testing}
\fancyhead[R]{
\includegraphics[width=4cm]{gfx/logo.pdf}
}
\fancyfoot[C]{\sffamily\fontsize{9pt}{9pt}\selectfont{\thepage/\pageref{LastPage}}}
\pagestyle{fancy}

\makeatletter

\newcommand{\spacedlowsmallcaps}[1]{\textbf{#1}}
\newcommand{\titleColor}{}
\makeatother

\usepackage{graphicx}
\usepackage[absolute]{textpos}
\usepackage[table]{xcolor}
\definecolor{lightgray}{gray}{0.9}

\begin{document}
\rowcolors{1}{}{lightgray}
\thispagestyle{empty}
\setlength{\TPHorizModule}{1mm}
\setlength{\TPVertModule}{1mm}

\begingroup
\fontsize{11pt}{11pt}\selectfont
\linespread{1.05}

\begin{textblock}{50}(0,0)
\noindent \includegraphics{gfx/00_cafoscari}
\noindent \end{textblock}
~
\begin{textblock}{143}(47,42)

\noindent {\LARGE{}Progetto per\\Ingegneria del Software \emph{(CT0090)}}\\
\vspace{30pt}

%\vspace{30pt}


\noindent %
\begin{minipage}[t]{1\columnwidth}%
\leavevmode
\titleColor
\raggedright
\Huge
\includegraphics[width=12cm]{gfx/logo.pdf}
\vspace{10pt}

WatchDoge:
Piano di Testing \\
\LARGE{Versione 1.0 - 15 novembre 2018}
\end{minipage}

\vspace{150pt}


\noindent {\large{}\spacedlowsmallcaps{{\large{}Redazione}}}{\large \par}
\begin{itemize}
\item   Dario Lazzaro
\item   Giovanni Scodeller
\item   Giulio Zausa
\item   Samuele Casarin
\item   Sandro Baccega
\end{itemize}

\vspace{20pt}


\noindent {\large{}\spacedlowsmallcaps{{\large{}Approvazione}}}{\large \par}
\begin{itemize}
\item   Dario Lazzaro
\item   Samuele Casarin
\item   Sandro Baccega
\end{itemize}
\end{textblock}
\endgroup
\pagebreak

%
%	inizio roba vera e propria
%

\tableofcontents

\pagebreak
  
\section{Introduzione}

Il seguente documento contiene le informazioni relative alla
pianificazione del processo di testing per il progetto WatchDoge. Tale
processo ha lo scopo di verifica e validazione dei requisiti di progetto
e verrà eseguito dal team di progetto durante e in seguito alla fase di
sviluppo.

La struttura del documento è la seguente:
\begin{itemize}
\item \textbf{Strategia di Testing}: in questa sezione sono definite le procedure adottate per il processo di testing in ognuno degli elementi considerati;
\item \textbf{Tracciabilità dei Requisiti}: in questa sezione sono definite le associazioni tra ciascuna specifica dei requisiti di progetto (vedi \emph{Documento di spefica dei requisiti}, sezione 3.5) e il metodo di testing e la condizione per cui il test risulta essere soddisfatto;
\item \textbf{Elementi di Testing}: in questa sezione sono riportati gli elementi soggetti al processo di testing;
\item \textbf{Schedule di Testing}: in questa sezione è visualizzato lo schema di allocazione temporale delle risorse di progetto impiegate per il processo di testing;
\item \textbf{Procedure di Registrazione dei Test}: in questa sezione sono definite le modalità con cui i risultati del processo ditesting vengono raccolti ed archiviati;
\item \textbf{Requisiti Hardware e Software}: in questa sezione sono definiti i requisiti in termini di hardware e software che devono essere soddisfatti per l'esecuzione del processo di testing;
\item \textbf{Vincoli}: in questa sezione sono definiti i vincoli da rispettare al fine che lo stesso processo di testing possa essere considerato attendibile.
\end{itemize}

Per una completa comprensione del presente documento e del progetto
WatchDoge, fare riferimento anche ai documenti \emph{Piano di Progetto}
e \emph{Documento di Specifica dei Requisiti}.

\section{Strategia di Testing}

Seguiremo due percorsi di testing separati: uno per il Robot ed uno per
l'applicazione Android. Adotteremo una strategia di testing basata
sull'approccio \textbf{Bottom-Up}: si andranno dapprima a collaudare i moduli di
basso livello presenti nella nostra gerarchia di astrazione, crescendo poi
di livello proseguendo con lo sviluppo. In questo modo, la verifica dei
moduli implementati partirà dalle unità elementari fino ad arrivare a
moduli sempre di più complessi.

Quando le due macro-componenti (Robot ed applicazione) avranno raggiunto
singolarmente una discreta stabilità, uniremo i due percorsi di testing,
in modo tale da poter effettuare test basati sull'interazione tra di
esse.

Suddivideremo il processo di testing in tre fasi:

\begin{enumerate}
\item
  Inizialmente, sarà data priorità al corretto funzionamento delle
  specifiche del \textbf{Robot}: essendo gran parte dell'applicazione dipendente
  a comportamenti basilari della macchina, è necessario ottenere un
  certo grado di sicurezza sul corretto funzionamento di quest'ultima.
\item
  In seguito, potremo cominciare la fase di testing relativa
  all'applicazione \textbf{Android}; non potendo avere a disposizione in questa
  fase l'implementazione completa delle funzionalità utili per
  l'interazione con il Robot, ricorreremo anche all'utilizzo di dummy
  objects. Per queste due prime fasi, adotteremo per lo più la strategia
  di testing \textbf{White-Box}, per avere così una certa garanzia di correttezza
  della parte software sviluppata indipendentemente.
\item
  Una volta raggiunto un buon punto di stabilità per lo sviluppo singolo
  del Robot e dell'applicazione, sarà possibile e necessario verificare
  che \textbf{l'interazione tra le due macro-componenti} sia quella desiderata e,
  probabilmente, la stessa ottenuta dai test eseguiti con dummy objects.
  In quest'ultima fase si adotterà maggiormente la strategia di testing
  \textbf{Black-Box}, ovvero simulando le possibili azioni che un utente finale
  potrà compiere.
\end{enumerate}

Altre strategie di testing che adotteremo sono le seguenti:
\begin{itemize}
\item \textbf{Incremental Testing: }i test vengono effettuati all'aggiunta di nuove unità, per evitare la regressione sulle unità precedentemente sviluppate;
\item \textbf{Stress Testing: }i test vengono effettuati per verificare che il
sistema sopporti il carico massimo definito in fase di progettazione
(tipo di test necessario per sistemi distribuiti);
\item \textbf{Analisi Statica:}
verificare la corrispondenza tra un sistema software e la sua specifica
senza eseguire il codice, utilizzando software che ne verifichino il contenuto in maniera statica.
\end{itemize}

\section{Tracciabilità dei
requisiti}

Gli identificatori utilizzati per la creazione della tabella sono
definiti nel Documento di Specifica dei Requisiti, sezione 3.5.

Tutti i test cases descritti di seguito devono essere riproducibili e
rispettare i risultati previsti.

\begin{figure}[htbp]
\centering
\begin{tabular}{| c | p{4cm} p{4cm} p{4cm} |}
\hline
\textbf{ID Requisito} & \textbf{Nome Requisito} & \textbf{Test Alpha} & \textbf{Test Beta} \\ \hline
FRS-01 & Associazione con codice Pin & Si verifica quando il Doge riconosce il Pin inviato dal telefono. & Si verifica quando il Doge ed il telefono si associano
correttamente. \\ \hline
FRS-02 & Rilevazione di un intruso & Il Doge riesce a vedere l'ambiente attraverso la telecamera. & Il Doge riesce a rilevare un oggetto in movimento. \\ \hline
FRS-03 & Notificazione acustica dell'intrusione & Il Doge emette un suono acustico. & Alla rilevazione di un movimento emette un suono acustico \\ \hline
FRS-04 & Disattivazione antifurto & Il Doge si spegne attraverso un segnale. & Il Doge si spegne tramite una sequenza di tasti. \\ \hline
FRS-05 & Primo avvio & L'operazione è verificata se al tap l'applicazione si avvia
correttamente. & L'operazione è verificata se gli utenti riescono ad aprire
l'applicazione correttamente. \\ \hline
FRS-06 & Utilizzo menù & Controllo del menù, apertura e primo utilizzo. & Utilizzo stressante dell'applicazione. \\ \hline
FRS-07 & Programmazione Doge & Controllo delle funzionalità programmabili. & Inserimento di tutte le possibili combinazioni programmabili. \\ \hline
FRS-08 & Sistema di notifica & L'applicazione crea notifiche. & Le notifiche vengono create con informazioni inviate dal robot. \\ \hline
FRS-09 & Spostamento manuale Doge & Il Doge riesce a muoversi su un asse. & Il Doge riesce a muoversi attraverso dei segnali. \\
\hline
\end{tabular}
\end{figure}

\section{Elementi di Testing}

Gli elementi sottoposti ai vari test durante il periodo di debugging
saranno principalmente tutte le componenti presenti nel Robot e
nell'applicazione.

Prove specifiche finalizzate a scoprire e risolvere eventuali
problematiche che potrebbero presentarsi nel Robot riguarderanno:
\begin{itemize}
    \item I singoli componenti forniti nel set \textbf{Lego Mindstorm EV3};
    \item Il \textbf{firmware} del Robot;
    \item Le unità, a partire da quelle elementari alle più complesse,
dell'applicazione \textbf{Android}.
\end{itemize}

\section{Schedule di Testing}

\begin{longtable}[]{@{}lllll@{}}
\toprule
Nome Requisito & Tipo test & Tester & Data & Esito\tabularnewline
\midrule
\endhead
Associazione con codice Pin & Alfa & & &\tabularnewline
Rilevazione di un intruso & Alfa & & &\tabularnewline
Notificazione acustica dell'intrusione & Alfa & & &\tabularnewline
Disattivazione antifurto & Alfa & & &\tabularnewline
Primo avvio & Alfa & & &\tabularnewline
Utilizzo menù & Alfa & & &\tabularnewline
Programmazione Doge & Alfa & & &\tabularnewline
Sistema di notifica & Alfa & & &\tabularnewline
Spostamento manuale Doge & Alfa & & &\tabularnewline
Associazione con codice Pin & Beta & & &\tabularnewline
Rilevazione di un intruso & Beta & & &\tabularnewline
Notificazione acustica dell'intrusione & Beta & & &\tabularnewline
Disattivazione antifurto & Beta & & &\tabularnewline
Primo avvio & Beta & & &\tabularnewline
Utilizzo menù & Beta & & &\tabularnewline
Programmazione Doge & Beta & & &\tabularnewline
Sistema di notifica & Beta & & &\tabularnewline
Spostamento manuale Doge & Beta & & &\tabularnewline
\bottomrule
\end{longtable}

\section{Procedure di Registrazione dei
Test}

Per eseguire ed archiviare il risultato dei test ci avvaleremo di
\emph{Travis CI}, uno strumento software gratuito messo a disposizione
su \emph{Github} per il testing del software.

Dopo la configurazione iniziale, \emph{Travis CI} verificherà ogni commit della repository remota del progetto su Github, eseguendo dei test automatici e tenendo traccia dell'esito automaticamente.

Sarà possibile accedere agli esiti dei test effettuati da \emph{Travis
CI} tramite un'apposita sezione su Github.

Per tutto il codice per cui \emph{Travis CI} non è configurato,
sfrutteremo il modello definito nella sezione 5 \emph{Schedule di
Testing} per documentare il processo di testing: descrizione della
funzionalità, nome del tester, data di testing ed esito.

\section{Requisiti Hardware e
Software}

Per eseguire i vari test, useremo questi strumenti:

\begin{itemize}
\item
  \textbf{Hardware:}

  \begin{itemize}
  \item
    Smartphone personali dei componenti del gruppo per verificare
    l'applicazione Android
  \item
    Robot EV3 per verificare il programma
  \end{itemize}
\item
  \textbf{Software:}

  \begin{itemize}
  \item
    Emulatore di smartphone Android (disponibile nativamente su Android
    Studio)
  \item
    \emph{Travis CI} (disponibile su Github)
  \end{itemize}
\end{itemize}

\section{Vincoli}

Il processo di testing verrà effettuato durante e in seguito l'attività
di sviluppo del progetto e dovrà terminare entro 5 giorni prima della
data di scadenza (entro il 26 gennaio 2019), al fine di verificare e
validare tutti i requisiti di progetto entro i limiti di tempo e
riducendo il più possibile la probabilità di errori durante la fase di
produzione.


\end{document}
